% Generated draft: requires mathematical and source review.
\documentclass[11pt]{article}
\usepackage[margin=25mm]{geometry}
\usepackage{fontspec}
\setmainfont{Latin Modern Roman}
\usepackage{amsmath,amssymb,mathtools,unicode-math}
\setmathfont{Latin Modern Math}
\usepackage{enumitem,needspace,xurl,hyperref}
\setlength{\parindent}{0pt}
\setlength{\parskip}{5pt}
\setlength{\emergencystretch}{3em}
\newcommand{\N}{\mathbb{N}}
\newcommand{\Z}{\mathbb{Z}}
\newcommand{\Q}{\mathbb{Q}}
\newcommand{\R}{\mathbb{R}}
\newcommand{\C}{\mathbb{C}}
\newcommand{\F}{\mathbb{F}}
\newcommand{\A}{\mathbb{A}}
\newcommand{\PP}{\mathbb{P}}
\newcommand{\E}{\mathbb{E}}
\newcommand{\eps}{\varepsilon}
\newcommand{\dd}{\,\mathrm{d}}
\begin{document}
\section*{UnsolvedMath ID: 30002705}
\textbf{Source problem code:} catalogue entry
\subsection{Definitions and mathematical statement}
Let \((M,g)\) be a Riemannian manifold, let \(o\in M\), and let \(W=(W_s)_{s\geq 0}\) be Brownian motion on \(M\) with \(W_0=o\). Let \(N\subset M\) be a hypersurface which locally separates \(M\) into sides denoted by \(M^+\) and \(M^-\). For any specified measurable positive side \(M^+\subseteq M\) and any deterministic time \(t\geq 0\) through which the process is defined, its occupation time in that side is
\[
T_t:=\int_0^t \mathbf{1}_{M^+}(W_s)\,\mathrm{d}s,
\qquad 0\leq T_t\leq t.
\]
The requested problem is to describe this occupation time. A distributional reading asks for the probability law \(\mathcal{L}_o(T_t)\) for the given geometric data and arbitrary such \(t\); the source does not specify the required form of a description.

The definition above is conditional on having a specified positive side wherever the path travels before time \(t\). Local separation by \(N\) alone does not specify such a global set. If the sides exist only in a neighborhood \(U\), one could instead define
\[
\tau_U:=\inf\{s\geq 0:W_s\notin U\},
\qquad
T_t^U:=\int_0^{t\wedge\tau_U}\mathbf{1}_{M^+\cap U}(W_s)\,\mathrm{d}s,
\]
provided a side in \(U\) is specified. This stopped quantity is a possible clarification, not an identification of the source's intended random variable. The source does not state that \(o\in N\), fix the normalization of the Brownian generator, specify boundary behavior or lifetime conventions, or supply further regularity hypotheses. These qualifications remain unresolved.
\subsection{Short English statement}
Describe how long Brownian motion started at \(o\) spends on the positive side of a separating hypersurface before time \(t\), naturally by determining its probability distribution. The source leaves the global meaning of the positive side and several process conventions unspecified. Its supplied partially solved status concerns a dated short-time assessment and requires independent review.
\subsection{Sources}
\begin{itemize}
\item Local UnsolvedMath registry, \texttt{data/unsolvedmath/problems.json}, permanent ID \(30002705\), title \emph{Occupation Times of Manifold Brownian Motion}, exact source code \texttt{catalogue entry}. The supplied record assigns proposed year 2014 and names no individual proposer. It attributes the problem to Oberwolfach workshop open-problem discussions.
\item Supplied report citation: \emph{Dirichlet Form Theory and its Applications (2015)}, Oberwolfach Reports, pp.\ 2667--2756, DOI \texttt{10.4171/owr/2014/48}. \url{https://doi.org/10.4171/owr/2014/48}. The supplied literature list also labels this reference \emph{Oberwolfach Report 48/2014}.
\item Supplied literature reference: Hsu--Ouyang, \emph{Geometric Deviation from Levys Law}. \url{https://arxiv.org/abs/2004.10039}.
\item Literature assessment dated exactly \texttt{2026-08-21}: the supplied curation labels the problem partially solved and reports that Hsu and Ouyang obtain a rigorous short-time expansion for the occupation time, including the arcsine term and the mean-curvature/local-time correction, while a complete distributional description for arbitrary time is not supplied. This is supplied dated curation derived from the OWR record, not an independently verified literature assessment; its status and cited claims require independent verification.
\end{itemize}
\end{document}
